\chapter{La chaîne de traçabilité}\label{ch:trace}

\begin{objectifs}[Objectifs]
\begin{itemize}
  \item suivre un objectif jusqu'à son indicateur de production, et inversement ;
  \item repérer à quelle étape une rupture de la chaîne se détecte ;
  \item répondre à la question : « pourquoi ce code existe-t-il ? »
\end{itemize}
\end{objectifs}

Tout ce manuel converge vers une idée : \emph{chaque ligne de code doit pouvoir justifier son existence}. Voici, pour \sika{}, le chemin complet
d'un objectif, de la première réunion à la mesure en production.

\begin{figure}[H]
\centering
\resizebox{\linewidth}{!}{%
\begin{tikzpicture}[
  n/.style={draw,rounded corners=3pt,align=center,font=\sffamily\small,minimum width=3.3cm,minimum height=13mm},
  c0/.style={n,fill=accent!8,draw=accent},
  c1/.style={n,fill=ochre!12,draw=ochre},
  c2/.style={n,fill=leaf!12,draw=leaf},
  ar/.style={-{Stealth},thick,grey},
  l/.style={font=\sffamily\scriptsize,text=grey,align=center}]
  \node[c0] (o) at (0,0) {\textbf{O1}\\cotisations contestées};
  \node[c0] (f) at (4.4,0) {\textbf{Fonctionnalité}\\Cotiser par MoMo};
  \node[c1] (s) at (8.8,0) {\textbf{SPEC-001}\\règle R2};
  \node[c1] (g) at (13.2,0) {\textbf{.feature}\\scénario \q{Confirmation}};
  \node[c2] (a) at (13.2,-2.6) {\textbf{openapi.yaml}\\submitContribution};
  \node[c2] (t) at (8.8,-2.6) {\textbf{Test}\\\texttt{test\_confirm...}};
  \node[c2] (r) at (4.4,-2.6) {\textbf{Recette}\\SC-02};
  \node[c0] (m) at (0,-2.6) {\textbf{Production}\\\% contestées à J+30};
  \draw[ar] (o)--(f); \draw[ar] (f)--(s); \draw[ar] (s)--(g); \draw[ar] (g)--(a); \draw[ar] (a)--(t); \draw[ar] (t)--(r); \draw[ar] (r)--(m);
  \draw[ar,dashed,accent] (m.north)--(o.south) node[midway,right,l]{boucle de retour};
\end{tikzpicture}}
\caption{La chaîne de traçabilité d'un objectif de \sika{} : chaque maillon est contrôlé par un artefact ou un garde-fou.}
\end{figure}

\section{Où une rupture se détecte}
\noindent\begin{tabularx}{\linewidth}{>{\raggedright\arraybackslash}p{4.6cm}X}
\toprule
\textbf{Rupture} & \textbf{Détectée par} \\ \midrule
Fonctionnalité sans objectif & colonne « Objectif » de l'impact map, \skill{forge-cadrage} \\
Terme utilisé sous deux noms & \cmd{check-glossary-terms.sh} (FF-05) \\
Code sans spec validée & hook \emph{pre-commit}, \cmd{check-spec-validated.sh} \\
Opération d'API sans spec & \cmd{check-openapi-first.sh} \\
Décision sans trace & \cmd{check-adr-present.sh}, ADR \\
Contexte qui en importe un autre & FF-01 en CI \\
Travail qui déborde du cycle court & \cmd{check-rapide-eligible.sh} \\
Production sans mesure & plan de monitoring, revues J+7 et J+30 \\
\bottomrule
\end{tabularx}

\section{Votre carnet de bord, version finale}
\begin{vousjouez}[À vous de jouer : votre propre chaîne]
Choisissez \emph{une} fonctionnalité de votre projet. Écrivez, à la main, les huit maillons de la figure ci-dessus avec les identifiants réels : objectif,
spec, règle, scénario, opération, test, scénario de recette, indicateur. S'il vous manque un maillon, vous venez de trouver la prochaine tâche du projet.
\end{vousjouez}

\begin{retenir}[À retenir]
\begin{itemize}
  \item Chaque ligne de code remonte à une règle, une spec, une fonctionnalité, un objectif.
  \item Les garde-fous contrôlent les maillons mécaniques ; la relecture humaine contrôle le sens.
  \item La boucle se ferme : la mesure en production nourrit le cadrage du lot suivant.
\end{itemize}
\end{retenir}
